% Generated by tools/build_new_dists_anthology_sections.py; do not edit by hand.

\section{Solvable and finite-subgroup families}\label{proof:new-dist:solvable_finite}

\resultbox{\textbf{Canonical testing artifacts.}\par\smallskip Exact backend: \path{new_dists/solvable_finite_subgroup_sigma_mgfs/verification.py}.\par Regression: \path{new_dists/solvable_finite_subgroup_sigma_mgfs/test_verification.py}.\par Marimo: \path{marimo_notebooks/solvable_finite.py}.}

\subsection*{Scope and outcome}

A group name alone does not determine a $\sigma$-MGF: the series depends on a
specified complex representation.  We prove the formulas for two distinct
finite-lamplighter representations, split metacyclic induced modules, finite
affine and Frobenius permutation actions, unitriangular, Borel, and parabolic
actions, finite defining subgroups of $SU(3)$, $SU(4)$, $Sp(4)$, and $SO(4)$,
and symplectic-reflection wreath products.  The accompanying marimo notebook
constructs every representative matrix group and compares direct Reynolds
averages with weighted-cycle, fixed-power, class-spectrum, or wreath-recursion
formulas.  Dimensions range from $2$ to $27$; the audit performs $7{,}580$
fixed-power checks and $70$ coefficient comparisons, all passing.  The finite
theorems, computational checks, and asymptotic statements are kept logically
separate.

\subsection{The representation-dependent target}

Let $\rho:G\to GL(W)$ be a finite complex representation.  For a partition
or composition $\tau=(\tau_1,\ldots,\tau_s)$, put
\[
 W_\tau=\bigotimes_{j=1}^s\Sym^{\tau_j}W,
 \qquad
 a_\tau(G,W)=\dim W_\tau^G.
\]
The normalization used here is
\[
 \MGF_{G,W}(\mathbf t)=\sum_\tau a_\tau(G,W)m_\tau(\mathbf t).
\]

\begin{theorem}[generalized Reynolds--Molien identity]\label{nd:solvable_finite:thm:universal}
For every finite complex representation,
\begin{align}
 \MGF_{G,W}(\mathbf t)
 &=\frac1{|G|}\sum_{g\in G}\prod_{i\geq1}
       \det(I-t_i\rho(g))^{-1},\label{nd:solvable_finite:eq:molien}\\
 &=\frac1{|G|}\sum_{g\in G}
 \exp\!\left(\sum_{r\geq1}\frac{p_r(\mathbf t)}r\chi_W(g^r)\right),
 \label{nd:solvable_finite:eq:power-trace}\\
 [m_\tau]\MGF_{G,W}&=a_\tau(G,W).\label{nd:solvable_finite:eq:coefficient}
\end{align}
\end{theorem}

\begin{proof}
The Reynolds operator $|G|^{-1}\sum_g\rho_{W_\tau}(g)$ is an idempotent with
image $W_\tau^G$, so its rank equals its trace.  The symmetric-power trace
identity
\[
 \sum_{q\geq0}\operatorname{Tr}(\Sym^q A)t^q=\det(I-tA)^{-1}
\]
gives \eqref{nd:solvable_finite:eq:molien} coefficient by coefficient.  Finally,
$-\log\det(I-tA)=\sum_{r\geq1}t^r\operatorname{Tr}(A^r)/r$ gives
\eqref{nd:solvable_finite:eq:power-trace}.
\end{proof}

If $G$ acts on a finite set $X$, we always mean the honest complex
permutation module $W=\C[X]$.  Write
\[
 F_r(g)=|\Fix_X(g^r)|,
 \qquad
 c_d(g)=\#\{\text{$d$-cycles of $g$ on $X$}\}.
\]

\begin{corollary}[permutation form]\label{nd:solvable_finite:cor:permutation}
For every finite permutation action,
\begin{align}
 F_r(g)&=\sum_{d\mid r}d c_d(g),&
 c_d(g)&=\frac1d\sum_{r\mid d}\mu(d/r)F_r(g),\label{nd:solvable_finite:eq:mobius}\\
 \MGF_{G,X}&=\frac1{|G|}\sum_{g\in G}
 \prod_{i,d\geq1}(1-t_i^d)^{-c_d(g)},\label{nd:solvable_finite:eq:perm-molien}\\
 [m_\tau]\MGF_{G,X}
 &=\#\left(G\backslash\prod_j\Sym^{\tau_j}X\right).\label{nd:solvable_finite:eq:orbit}
\end{align}
In particular, $[m_{(1^s)}]\MGF=\E F_1(g)^s=\#(G\backslash X^s)$.
\end{corollary}

\begin{proof}
A point in a $d$-cycle is fixed by $g^r$ exactly when $d\mid r$; M\"obius
inversion gives \eqref{nd:solvable_finite:eq:mobius}.  A $d$-cycle has determinant factor
$1-t^d$, which proves \eqref{nd:solvable_finite:eq:perm-molien}.  Monomials in
$\Sym^q\C[X]$ are indexed by $q$-element multisets, and orbit sums form a
basis for the invariant space, proving \eqref{nd:solvable_finite:eq:orbit}.
\end{proof}

\subsection{Finite lamplighter groups}

\subsubsection{Natural monomial representation}

Put $L_{r,n}=\mu_r^n\rtimes C_n$, where the generator $P$ cyclically shifts
the coordinates, and use the $n$-dimensional representation
\[
 (z_1,\ldots,z_n;P^s)\longmapsto
 \operatorname{diag}(z_1,\ldots,z_n)P^s.
\]
Define
\[
 A_{r,\ell}(\mathbf t)=\frac1r\sum_{\zeta^r=1}
 \prod_{i\geq1}(1-\zeta t_i^\ell)^{-1}.
\]
Root-of-unity filtering also gives
\[
 A_{r,\ell}=\sum_{q\geq0}h_{qr}(t_1^\ell,t_2^\ell,\ldots).
\]

\begin{theorem}[lamplighter monomial formula]\label{nd:solvable_finite:thm:lamp-monomial}
\[
 \boxed{\MGF^{\mathrm{mon}}_{L_{r,n}}
 =\frac1n\sum_{\ell\mid n}\varphi(\ell)
 A_{r,\ell}^{\,n/\ell}.}
\]
\end{theorem}

\begin{proof}
The shift $P^s$ has $d=\gcd(n,s)$ cycles, each of length
$\ell=n/d$.  On one coordinate cycle $C$, weighted permutation elimination
gives
\[
 \det(I-t\operatorname{diag}(z)P^s\mid_C)
 =1-\left(\prod_{j\in C}z_j\right)t^\ell.
\]
When the labels are averaged independently, every cycle product is an
independent uniform element of $\mu_r$.  Thus a shift of cycle length $\ell$
contributes $A_{r,\ell}^{n/\ell}$.  Exactly $\varphi(\ell)$ shifts have this
cycle length.
\end{proof}

There is no raw coefficientwise limit with $r$ fixed and $n\to\infty$.  The
identity-shift sector has weight $1/n$ and contributes $A_{r,1}^n/n$.
Choosing the $h_r$ term from two factors gives
\[
 [m_{(r,r)}]\MGF^{\mathrm{mon}}_{L_{r,n}}
 \geq \frac1n\binom n2[m_{(r,r)}]h_r^2
 =\frac{(r+1)(n-1)}2.
\]

\subsubsection{Affine action on lamp configurations}

The same abstract group acts on $X=C_r^n$ by $x\mapsto b+P^sx$.  For
$g=(b,P^s)$ set $b_k=(I+P^s+\cdots+P^{(k-1)s})b$.

\begin{proposition}[lamp fixed powers and pair coefficient]
\begin{align}
 F_k(g)&=\begin{cases}
 r^{\gcd(n,ks)},&b_k\in\im(I-P^{ks}),\\
 0,&\text{otherwise},
 \end{cases}\label{nd:solvable_finite:eq:lamp-affine-fixed}\\
 [m_{(1,1)}]\MGF^{\mathrm{aff}}_{L_{r,n}}
 &=\boxed{\frac1n\sum_{d\mid n}\varphi(d)r^{n/d}}.\label{nd:solvable_finite:eq:necklace}
\end{align}
\end{proposition}

\begin{proof}
Iteration gives $g^k(x)=b_k+P^{ks}x$.  Hence fixed points solve
$(I-P^{ks})x=b_k$.  A soluble fiber has the size of the kernel, namely
$r^{\gcd(n,ks)}$, proving \eqref{nd:solvable_finite:eq:lamp-affine-fixed}.  For ordered pairs,
translate the first entry to $0$.  The remaining difference is acted on only
by cyclic shifts.  Burnside's lemma gives the necklace count
\eqref{nd:solvable_finite:eq:necklace}.
\end{proof}

For fixed $r\geq2$, the necklace count is asymptotic to $r^n/n$; thus this
representation has exponential degree-two growth.  This does not contradict
Theorem~\ref{nd:solvable_finite:thm:lamp-monomial}: the two series use different representations.

\subsection{Split metacyclic induced representations}

Let
\[
 M(m,n,u)=\langle a,b:a^m=b^n=1,\ bab^{-1}=a^u\rangle,
 \qquad u^n\equiv1\pmod m,
\]
and suppose the chosen character has a full orbit of length $n$.  With
$\zeta=e^{2\pi i/m}$, define
\[
 \rho(a)e_j=\zeta^{u^{-j}}e_j,
 \qquad \rho(b)e_j=e_{j+1}.
\]
For $g=a^xb^y$, let $\mathcal C_y$ be the cycles of addition by $y$ on
$\mathbb Z/n\mathbb Z$, and put
\[
 \alpha_C(x,y)=\zeta^{x\sum_{j\in C}u^{-j}}.
\]

\begin{theorem}[metacyclic weighted-cycle formula]\label{nd:solvable_finite:thm:metacyclic}
\[
 \boxed{\MGF_{M(m,n,u),\rho}
 =\frac1{mn}\sum_{x=0}^{m-1}\sum_{y=0}^{n-1}
 \prod_{i\geq1}\prod_{C\in\mathcal C_y}
 (1-\alpha_C(x,y)t_i^{|C|})^{-1}.}
\]
\end{theorem}

\begin{proof}
The relations hold because shifting the diagonal weights sends
$u^{-j}$ to $u^{-(j-1)}=u\,u^{-j}$.  On a cycle $C$ of the permutation
$j\mapsto j+y$, the product of the diagonal labels is $\alpha_C(x,y)$.
Weighted permutation elimination gives the characteristic factor
$1-\alpha_C(x,y)t^{|C|}$.  Multiplication over cycles and Reynolds averaging
prove the result.
\end{proof}

If the character orbit has length $d<n$, the same proof applies to its
$d$-dimensional induced module.  A nonsplit relation $b^n=a^c$ inserts the
corresponding scalar in cycles crossing that relation; it must not be silently
omitted.  The arithmetic quantities $\gcd(n,y)$ and
$\sum_{j\in C}u^{-j}\pmod m$ prevent a group-independent asymptotic law.

\subsection{General affine and Frobenius permutation actions}

\subsubsection{Affine groups}

Let $V=\F_q^d$, $H\leq GL(V)$, and let $V\rtimes H$ act naturally on $V$.
For $g=(b,h)$ set $b_r=(I+h+\cdots+h^{r-1})b$.

\begin{theorem}[affine fixed powers and moments]\label{nd:solvable_finite:thm:affine}
\begin{align}
 F_r(b,h)&=\begin{cases}
 q^{\dim\ker(I-h^r)},&b_r\in\im(I-h^r),\\
 0,&\text{otherwise},
 \end{cases}\label{nd:solvable_finite:eq:affine-fixed}\\
 [m_{(1^s)}]\MGF_{V\rtimes H,V}
 &=\#(H\backslash V^{s-1}).\label{nd:solvable_finite:eq:affine-moment}
\end{align}
For $H=GL_d(q)$ this specializes to
\begin{equation}
 \boxed{[m_{(1^s)}]\MGF
 =\sum_{j=0}^{\min(d,s-1)}\qbinom{s-1}{j}.}\label{nd:solvable_finite:eq:gaussian}
\end{equation}
\end{theorem}

\begin{proof}
Iteration gives $g^r(x)=b_r+h^rx$, and solving the affine linear equation
gives \eqref{nd:solvable_finite:eq:affine-fixed}.  Translation sends the first point of an
$s$-tuple to zero and leaves the diagonal $H$-action on its $s-1$
differences, proving \eqref{nd:solvable_finite:eq:affine-moment}.  Regard an element of
$V^{s-1}$ as a $d$-by-$(s-1)$ matrix.  Two matrices lie in the same
$GL_d(q)$ orbit under left multiplication exactly when they have the same row
space.  The possible row spaces are the $j$-subspaces of $\F_q^{s-1}$ with
$j\leq d$, proving \eqref{nd:solvable_finite:eq:gaussian}.
\end{proof}

For fixed $q$ and fixed $s$, the right side stabilizes once $d\geq s-1$.
For general $H_d$, coefficientwise convergence is equivalent to stabilization
of the $H_d$-orbit counts on every fixed number of vectors.

\subsubsection{Finite Frobenius groups}

Let $G=K\rtimes H$ be a Frobenius group in its natural action on
$G/H\simeq K$, and put $N=|K|$.  Define
\[
 E_d=\prod_i(1-t_i^d)^{-N/d},\qquad
 R_d=\prod_i(1-t_i)^{-1}(1-t_i^d)^{-(N-1)/d}.
\]

\begin{theorem}[Frobenius cycle inventory]\label{nd:solvable_finite:thm:frobenius}
\begin{align}
 \MGF_{G,K}
 &=\boxed{\frac1{N|H|}\left[E_1+
 \sum_{1\ne k\in K}E_{\operatorname{ord}(k)}+
 N\sum_{1\ne h\in H}R_{\operatorname{ord}(h)}\right]},\label{nd:solvable_finite:eq:frob-molien}\\
 [m_{(1,1)}]\MGF_{G,K}
 &=\boxed{1+\frac{N-1}{|H|}}.\label{nd:solvable_finite:eq:frob-pair}
\end{align}
\end{theorem}

\begin{proof}
A nonidentity $k\in K$ is a fixed-point-free translation; if its order is
$d$, its cycles have type $d^{N/d}$.  Every element outside $K$ lies in a
unique conjugate of a complement and fixes one point.  Its nontrivial powers
fix the same point, so an element of order $d$ has cycle type
$1^1d^{(N-1)/d}$.  The $N$ complements have disjoint nonidentity parts,
which proves \eqref{nd:solvable_finite:eq:frob-molien}.  After translating the first member of
an ordered pair to $1\in K$, the complement acts freely on $K\setminus\{1\}$.
Its orbit count proves \eqref{nd:solvable_finite:eq:frob-pair}.
\end{proof}

With $|H|$ fixed and $N\to\infty$, the pair coefficient grows linearly.
Sharper transitivity can keep degree two bounded, but higher orbit parameters
still require separate control.

\subsection{Unitriangular, Borel, and parabolic actions}

Finite-field matrices are not complex matrices until a complex
representation is chosen.  Here the choices are permutation modules on
$V_n=\F_q^n$, $V_n^\times$, projective points, and flag varieties.

\begin{proposition}[vector and projective fixed powers]\label{nd:solvable_finite:prop:linear-fixed}
For $g\in GL_n(q)$,
\begin{align}
 F_r(g;V_n)&=q^{\dim\ker(g^r-I)},\label{nd:solvable_finite:eq:vector-fixed}\\
 F_r(g;V_n^\times)&=q^{\dim\ker(g^r-I)}-1,\label{nd:solvable_finite:eq:nonzero-fixed}\\
 F_r(g;\mathbb P^{n-1})
 &=\sum_{\lambda\in\F_q^\times}
 \frac{q^{\dim\ker(g^r-\lambda I)}-1}{q-1}.\label{nd:solvable_finite:eq:projective-fixed}
\end{align}
\end{proposition}

\begin{proof}
The first two statements count fixed vectors, with or without zero.  A
projective point is fixed exactly when its vectors form an eigenline.  The
nonzero vectors in distinct eigenspaces are disjoint, and every line has
$q-1$ representatives, giving \eqref{nd:solvable_finite:eq:projective-fixed}.
\end{proof}

For the upper-unitriangular and upper-triangular Borel groups, elementary row
operations give
\begin{align}
 \#(UT_n(q)\backslash V_n^\times)&=n(q-1),&
 \#(UT_n(q)\backslash\mathbb P^{n-1})&=n,\label{nd:solvable_finite:eq:ut-orbits}\\
 \#(B_n(q)\backslash V_n^\times)&=n,&
 \#(B_n(q)\backslash\mathbb P^{n-1})&=n.\label{nd:solvable_finite:eq:b-orbits}
\end{align}
For $UT_n(q)$ a nonzero vector is classified by the position and value of its
last nonzero coordinate; projectivization removes the value.  The diagonal
torus in $B_n(q)$ removes it before projectivization.  On all of $V_n$, add
one orbit for zero.  Consequently these rank-growing actions already diverge
in degree one.

Let $G$ be a finite reductive group and let $P,Q$ be parabolics with Weyl
subgroups $W_P,W_Q$.  Bruhat decomposition and the coset fixed-point count
give
\begin{align}
 [m_{(1)}]\MGF_{P,G/Q}&=|P\backslash G/Q|=|W_P\backslash W/W_Q|,
 \label{nd:solvable_finite:eq:parabolic-orbits}\\
 |\Fix_{G/Q}(p^r)|&=
 \frac{|C_G(p^r)|\,|(p^r)^G\cap Q|}{|Q|}.\label{nd:solvable_finite:eq:parabolic-fixed}
\end{align}
For $P=Q=B_n(q)$, \eqref{nd:solvable_finite:eq:parabolic-orbits} is $|W|=n!$.
Equation~\eqref{nd:solvable_finite:eq:parabolic-fixed} follows by counting the $x\in G$ for
which $x^{-1}p^rx\in Q$, then dividing by $|Q|$.

\subsection{Finite defining subgroups of unitary and symplectic groups}

\subsubsection{Class spectra and scalar selection}

For $G\leq U(d)$ let $\lambda_1(g),\ldots,\lambda_d(g)$ be the eigenvalues
of the defining matrix.

\begin{proposition}[finite unitary class formula]\label{nd:solvable_finite:prop:unitary-class}
\begin{align}
 \MGF_{G,\C^d}
 &=\boxed{\sum_{[g]}\frac1{|C_G(g)|}
 \prod_{i\geq1}\prod_{j=1}^d(1-\lambda_j(g)t_i)^{-1}},
 \label{nd:solvable_finite:eq:unitary-class}\\
 &=\sum_{[g]}\frac1{|C_G(g)|}
 \exp\!\left(\sum_{r\geq1}\frac{p_r}{r}\chi(g^r)\right).
 \label{nd:solvable_finite:eq:unitary-power}
\end{align}
Thus a character table determines the complete series only together with its
class power maps.  If $G$ contains a scalar $\zeta I$ of order $s$, then
\[
 [m_\tau]\MGF_G=0\quad\text{unless}\quad s\mid|\tau|.
 \label{nd:solvable_finite:eq:scalar-selection}
\]
\end{proposition}

\begin{proof}
Group the average in Theorem~\ref{nd:solvable_finite:thm:universal} by conjugacy classes to get
\eqref{nd:solvable_finite:eq:unitary-class}; the logarithmic determinant identity gives
\eqref{nd:solvable_finite:eq:unitary-power}.  The scalar acts on $W_\tau$ as
$\zeta^{|\tau|}$, so a fixed vector requires that scalar to be one.
\end{proof}

\subsubsection{General monomial theorem and the SU(3) type-C and type-D series}

Let $G=D\rtimes P\leq U(d)$ with $D$ diagonal and $P\leq S_d$.  For a
cycle $C$ of $\pi\in P$, write $z_C=\prod_{j\in C}z_j$.

\begin{theorem}[monomial class formula]\label{nd:solvable_finite:thm:monomial}
\[
 \boxed{\MGF_{D\rtimes P}
 =\frac1{|D||P|}\sum_{\pi\in P}\sum_{z\in D}
 \prod_{i\geq1}\prod_{C\in\Cyc(\pi)}
 (1-z_Ct_i^{|C|})^{-1}.}
\]
\end{theorem}

\begin{proof}
Weighted permutation elimination gives
$\det(I-tD(z)P_\pi)=\prod_C(1-z_Ct^{|C|})$.  Theorem~\ref{nd:solvable_finite:thm:universal}
then applies.  This includes the determinant-one monomial type-C and type-D
families in $SU(3)$.
\end{proof}

If $D_m$ is the $m$-torsion in a fixed compact diagonal torus $T$ normalized
by $P$, then in total degree at most $D$ only characters of bounded exponent
occur.  For all sufficiently large $m$, such a character is trivial on
$D_m$ exactly when it is trivial on $T$.  Hence
\[
 \MGF_{D_m\rtimes P}\longrightarrow\MGF_{T\rtimes P}
 \quad\text{coefficientwise}.
\]
For the full coordinate torus and the determinant-one torus, $m>D$ suffices.

\subsubsection{Two exceptional SU(3) class formulas}

Write
\[
 \mathcal D(\alpha,\beta,\gamma)
 =\prod_{i\geq1}\frac1{(1-\alpha t_i)(1-\beta t_i)(1-\gamma t_i)}.
\]
For the $3$-dimensional rotation representation of
$\Sigma(60)\cong A_5$,
\begin{align}
 \MGF_{\Sigma(60)}=\frac1{60}\bigl[&
 \mathcal D(1,1,1)+15\mathcal D(1,-1,-1)
 +20\mathcal D(1,\omega,\omega^2)\notag\\
 &+12\mathcal D(1,\zeta_5,\zeta_5^4)
 +12\mathcal D(1,\zeta_5^2,\zeta_5^3)\bigr].
 \label{nd:solvable_finite:eq:a5-su3}
\end{align}
For either conjugate $3$-dimensional representation of
$\Sigma(168)\cong PSL_2(7)$,
\begin{align}
 \MGF_{\Sigma(168)}=\frac1{168}\bigl[&
 \mathcal D(1,1,1)+21\mathcal D(1,-1,-1)
 +56\mathcal D(1,\omega,\omega^2)
 +42\mathcal D(1,i,-i)\notag\\
 &+24\mathcal D(\zeta_7,\zeta_7^2,\zeta_7^4)
 +24\mathcal D(\zeta_7^3,\zeta_7^5,\zeta_7^6)\bigr].
 \label{nd:solvable_finite:eq:psl-su3}
\end{align}
Both identities are immediate from Proposition~\ref{nd:solvable_finite:prop:unitary-class} and
the displayed class spectra.  Isolated exceptional groups have no intrinsic
rank asymptotic.  A sequence whose scalar-center order tends to infinity has
series tending coefficientwise to $1$ by \eqref{nd:solvable_finite:eq:scalar-selection}.

\subsubsection{SU(4), Sp(4), and SO(4)}

For $G\leq SU(4)$, Proposition~\ref{nd:solvable_finite:prop:unitary-class} with its defining
character $\chi_4$ is the complete answer.  For $G\leq Sp(4)$, the spectrum
is $\{\lambda,\lambda^{-1},\mu,\mu^{-1}\}$ and hence
\[
 \MGF_G=\sum_{[g]}\frac1{|C_G(g)|}\prod_i
 \frac1{(1-\lambda_gt_i)(1-\lambda_g^{-1}t_i)
 (1-\mu_gt_i)(1-\mu_g^{-1}t_i)}.
\]
If $-I\in G$, all odd-total-degree coefficients vanish.

For $SO(4)$ use $\operatorname{Spin}(4)\simeq SU(2)_L\times SU(2)_R$.
If $\widetilde G$ is the inverse image of $G$, the complexified vector
representation is the tensor product of the two defining planes, so
\[
 \boxed{\MGF_{G,\C^4}
 =\frac1{|\widetilde G|}\sum_{(a,b)\in\widetilde G}
 \exp\!\left(\sum_{r\geq1}\frac{p_r}{r}
 \chi_2(a^r)\chi_2(b^r)\right).}
 \label{nd:solvable_finite:eq:so4-spin}
\]
The kernel element $(-I,-I)$ acts trivially, so averaging over the spin cover
equals averaging over the represented image.

\subsection{Symplectic-reflection wreath products}

Let $K\leq SL_2(\C)$ act on its defining symplectic plane $E$, and let
$G_n=K\wr S_n$ act on $E^{\oplus n}$.  Define
\[
 A_\ell(\mathbf t)=\frac1{|K|}\sum_{k\in K}
 \prod_{i\geq1}\det(I-t_i^\ell k)^{-1}.
\]

\begin{theorem}[wreath exponential and stable law]\label{nd:solvable_finite:thm:wreath}
\begin{align}
 \sum_{n\geq0}u^n\MGF_{K\wr S_n,E^{\oplus n}}
 &=\boxed{\exp\!\left(\sum_{\ell\geq1}\frac{u^\ell}{\ell}A_\ell\right)},
 \label{nd:solvable_finite:eq:wreath-exponential}\\
 \lim_{n\to\infty}\MGF_{K\wr S_n,E^{\oplus n}}
 &=\boxed{\exp\!\left(\sum_{\ell\geq1}\frac{A_\ell-1}{\ell}\right)
 =\ExpS(\MGF_K-1)}.
 \label{nd:solvable_finite:eq:wreath-limit}
\end{align}
The limit is coefficientwise.
\end{theorem}

\begin{proof}
For a permutation cycle $C$ of length $\ell$, multiply its $K$-labels around
the cycle to obtain $k_C$.  Block elimination gives
$\det(I-tg_C)=\det(I-t^\ell k_C)$.  The cycle product is uniform in $K$, and
the exponential formula for labelled permutation cycles gives
\eqref{nd:solvable_finite:eq:wreath-exponential}.  Split $A_\ell=1+(A_\ell-1)$ to write the
right side as
\[
 \frac1{1-u}\exp\!\left(\sum_{\ell\geq1}
 \frac{u^\ell(A_\ell-1)}\ell\right).
\]
In a fixed positive $\mathbf t$-degree only finitely many $\ell$ occur, and
the second factor has bounded $u$-degree after that coefficient is extracted.
Multiplication by $(1-u)^{-1}$ makes the coefficients eventually constant,
proving \eqref{nd:solvable_finite:eq:wreath-limit}.
\end{proof}

This is a genuine marked compound-Poisson law, unlike the rare-sector
divergence of the lamplighter and fixed-complement Frobenius examples.

\subsection{Verification strategy and results}

\subsubsection{Independent computational routes}

The marimo notebook imports one pure verification module.  Every check uses
at least two of the following routes.

\begin{enumerate}
\item \textbf{Direct represented matrices.}  Dense complex matrices are used
for the monomial and defining representations.  A permutation matrix is
stored losslessly by its unique nonzero row in each column.
\item \textbf{Direct target quantity.}  For a matrix $A$, complete homogeneous
characters are computed from its eigenvalues, and
\[
 \frac1{|G|}\sum_g\prod_j h_{\tau_j}(\operatorname{Spec}\rho(g))
\]
is rounded only after its imaginary part and distance to the nearest integer
are checked.  Permutation coefficients are computed exactly from cycle
factors.  Ordered-pair and difference-tuple orbits are also enumerated
literally where applicable.
\item \textbf{Structural formula.}  Lamplighter and metacyclic cases use
root-filtered or label-product cycles; affine cases use image and kernel
tests; Frobenius cases use the predicted full cycle inventory; flag actions
use centralizers and class intersections; unitary cases use class spectra;
and wreath products use the coefficient recurrence obtained by differentiating
\eqref{nd:solvable_finite:eq:wreath-exponential}.
\item \textbf{Analytic spot check.}  The full determinant products are
evaluated at $(t_1,t_2)=(0.031,-0.047)$.  This checks all degrees at once and
is separate from the coefficient table.
\end{enumerate}

\subsubsection{Constructed groups and reasonable dimensions}

\begin{longtable}{@{}p{0.34\linewidth}rrrrp{0.15\linewidth}@{}}
\toprule
Representation & dim. & group/cover & fixed checks & coeff. checks & result\\
\midrule
\endhead
$L_{2,3},L_{3,2},L_{2,4}$, monomial & $2$--$4$ & $18$--$64$ & -- &15&\textcolor{teal}{\textbf{PASS}}\\
$L_{2,3},L_{2,4},L_{3,3}$, affine lamps &$8$--$27$&$24$--$81$&823&--&\textcolor{teal}{\textbf{PASS}}\\
$M(5,2,4),M(7,3,2)$, induced &$2$--$3$&$10,21$&relations&10&\textcolor{teal}{\textbf{PASS}}\\
$AGL(2,2),AGL(3,2)$, natural &$4,8$&$24,1344$&6634&5&\textcolor{teal}{\textbf{PASS}}\\
$\F_5\rtimes\F_5^\times$, $D_{14}$ &$5,7$&$20,14$&cycles&10&\textcolor{teal}{\textbf{PASS}}\\
$UT_3(2),B_2(3)$, vectors/projective &$7$--$8$&$8,12$&123&--&\textcolor{teal}{\textbf{PASS}}\\
$B_3(2)$ on complete flags &$21$&$8$&class powers&--&\textcolor{teal}{\textbf{PASS}}\\
$A_5,PSL_2(7),\Delta(12)$ in $SU(3)$ &$3$&$60,168,12$&relations&15&\textcolor{teal}{\textbf{PASS}}\\
$A_5$ in $SU(4)$ &$4$&$60$&classes&5&\textcolor{teal}{\textbf{PASS}}\\
$Q_8\wr S_2$ in $Sp(4)$ &$4$&$128$&center&5&\textcolor{teal}{\textbf{PASS}}\\
$(Q_8\times Q_8)/\{\pm(I,I)\}$ in $SO(4)$ &$4$&$64/32$&power traces&5&\textcolor{teal}{\textbf{PASS}}\\
\midrule
\textbf{Total}&&&\textbf{$7{,}580$}&\textbf{$70$}&\textcolor{teal}{\textbf{PASS}}\\
\bottomrule
\end{longtable}

For the $SU(3)$ exceptional cases the notebook constructs actual generators.
For $PSL_2(7)$, with $\eta=e^{2\pi i/7}$, it uses
\[
 A=\operatorname{diag}(\eta,\eta^2,\eta^4),\qquad
 B=\frac{i}{\sqrt7}
 \begin{pmatrix}
 \eta^2-\eta^5&\eta-\eta^6&\eta^4-\eta^3\\
 \eta-\eta^6&\eta^4-\eta^3&\eta^2-\eta^5\\
 \eta^4-\eta^3&\eta^2-\eta^5&\eta-\eta^6
 \end{pmatrix}.
\]
Closure produces $168$ distinct unitary determinant-one matrices with order
inventory $1,21,56,42,48$ for element orders $1,2,3,4,7$.  The maximum
unitarity and determinant errors are below $5.4\times10^{-14}$.

The $SO(4)$ test constructs all $64$ matrices $a\otimes b$ with
$a,b\in Q_8\leq SU(2)$.  Deduplication leaves $32$ represented matrices,
exactly detecting the kernel $\{(I,I),(-I,-I)\}$.  For powers $1\leq r\leq4$
it verifies
\[
 \operatorname{Tr}((a\otimes b)^r)
 =\operatorname{Tr}(a^r)\operatorname{Tr}(b^r)
\]
to machine precision, and the cover and image Molien values agree within
$6.7\times10^{-16}$.

\subsubsection{Representative coefficient values}

\begin{center}
\small
\begin{tabular}{@{}lrrrrr@{}}
\toprule
Representation & $m_1$ & $m_2$ & $m_{11}$ & $m_3$ & $m_{21}$\\
\midrule
$L_{2,3}$ monomial &0&1&1&0&0\\
$L_{3,2}$ monomial &0&0&0&1&1\\
$M(5,2,4)$ induced &0&1&1&0&0\\
$M(7,3,2)$ induced &0&0&0&1&1\\
$\F_5\rtimes\F_5^\times$ natural &1&2&2&3&5\\
$D_{14}$ natural &1&4&4&8&16\\
$A_5\leq SU(3)$ &0&1&1&0&0\\
$PSL_2(7)\leq SU(3)$ &0&0&0&0&0\\
$\Delta(12)\leq SU(3)$ &0&1&1&1&1\\
$A_5\leq SU(4)$ &0&1&1&1&1\\
$Q_8\wr S_2\leq Sp(4)$ &0&0&1&0&0\\
$(Q_8\times Q_8)/\pm\leq SO(4)$ &0&1&1&0&0\\
\bottomrule
\end{tabular}
\end{center}

The affine-lamp pair coefficients are $4,6,11$ for $(r,n)=(2,3),(2,4),(3,3)$,
matching literal pair orbits and the necklace formula.  For $AGL(2,2)$, the
moments $[m_{(1^s)}]$ for $s=2,3,4$ are $2,5,15$; for $AGL(3,2)$ the tested
values for $s=2,3$ are $2,5$.  In each case direct difference-tuple orbits,
affine Burnside averages, and Gaussian-binomial sums agree.

\subsection{Asymptotic summary and scope}

\begin{center}
\small
\begin{tabularx}{\linewidth}{@{}p{0.43\linewidth}X@{}}
\toprule
Family and representation & large-parameter behavior\\
\midrule
$C_r^n\rtimes C_n$, natural monomial
& $[m_{(r,r)}]\gg n$; no raw coefficientwise limit.\\
$C_r^n\rtimes C_n$, affine lamps
& $[m_{11}]\sim r^n/n$; no raw coefficientwise limit.\\
Split metacyclic induced modules
& arithmetic and subsequence dependent.\\
$V_d\rtimes GL_d(q)$, $q$ fixed
& every fixed $(1^s)$ moment stabilizes once $d\geq s-1$.\\
Frobenius groups with fixed complement
& $[m_{11}]=1+(|K|-1)/|H|$ grows linearly.\\
$UT_n(q),B_n(q)$ on nonzero/projective points
& $[m_1]$ grows linearly in $n$.\\
$B_n(q)$ on complete flags
& $[m_1]=n!$.\\
$SU(3)$ diagonal torsion series
& coefficientwise torus-normalizer limit as torsion becomes dense.\\
Exceptional fixed-dimensional unitary groups
& isolated exact class sums; no intrinsic rank limit.\\
$K\wr S_n\leq Sp(2n)$
& $\ExpS(\MGF_K-1)$, a genuine coefficientwise limit.\\
\bottomrule
\end{tabularx}
\end{center}

A full crystallographic group is generally infinite and noncompact; it has no
uniform probability measure and therefore no $\sigma$-MGF of the present
kind.  One must first choose a finite quotient, a compact quotient action, or
a nonuniform probability law.  A finite linear point group preserving a
complex lattice is covered by Proposition~\ref{nd:solvable_finite:prop:unitary-class}, or by
Theorem~\ref{nd:solvable_finite:thm:monomial} when it is monomial.

\resultbox{\textbf{Reproducibility.} Run
\texttt{python -m new\_dists.solvable\_finite\_subgroup\_sigma\_mgfs.verification}
for the noninteractive audit, or open
\path{marimo_notebooks/solvable_finite.py} with marimo for the
group-organized laboratory.  The source keeps the direct matrix
route and each family-specific formula in separate functions so that agreement
cannot result from comparing an expression with itself.}
